\chapter{\IfLanguageName{dutch}{Stand van zaken}{State of the Art}}%
\label{ch:stand-van-zaken}
In dit hoofdstuk wordt de theoretische achtergrond besproken die relevant is
voor de ontwikkeling van een webapplicatie voor realtime monitoring van
voertuigdata. Eerst wordt ingegaan op OBD-II en de manier waarop voertuigdata
via deze interface beschikbaar wordt gesteld. Vervolgens worden de
belangrijkste eigenschappen en beperkingen van het uitlezen van OBD-II-data
besproken.

Daarna wordt ingegaan op realtime dataverwerking binnen webapplicaties en
worden REST API's en WebSockets besproken als communicatiemethoden tussen
backend en frontend. Hierbij wordt aandacht besteed aan de eigenschappen van
beide technieken die relevant zijn voor het frequent doorsturen van
voertuigdata. Tot slot worden simulatie van voertuigdata en realtime
visualisatie kort besproken.

Deze theoretische achtergrond vormt de basis voor de verdere ontwikkeling van
de applicatie en voor de vergelijking tussen REST API en WebSockets die later
in deze bachelorproef wordt uitgevoerd.


\section{OBD-II}
On-Board Diagnostics (OBD) is een diagnosesysteem dat in voertuigen wordt
gebruikt om informatie over verschillende elektronische systemen beschikbaar
te stellen. De tweede generatie van deze technologie, OBD-II, introduceerde een
meer gestandaardiseerde manier om diagnostische informatie en voertuiggegevens
uit te lezen \autocite{mccord2011automotive}.

OBD-II werd oorspronkelijk voornamelijk ontwikkeld in het kader van
emissiegerelateerde diagnose. Door de standaardisering van de aansluiting en
een aantal diagnostische functies kan de interface echter ook worden gebruikt
om tijdens de werking van het voertuig verschillende parameters uit te lezen.
Hierdoor wordt OBD-II niet alleen gebruikt voor foutdiagnose, maar ook voor
toepassingen waarbij voertuiggegevens worden verzameld en geanalyseerd
\autocite{RIMPAS202055,Pmc2025}.

De OBD-II-aansluiting bevindt zich doorgaans in het interieur van het voertuig
en maakt het mogelijk om met behulp van een geschikte adapter communicatie met
de elektronische systemen van het voertuig tot stand te brengen. Een externe
applicatie kan via deze verbinding verzoeken naar het voertuig sturen en de
bijbehorende antwoorden verwerken.

OBD-II maakt daarbij gebruik van onderliggende communicatieprotocollen. In
moderne voertuigen wordt hiervoor vaak het Controller Area Network (CAN)
gebruikt. CAN is een communicatieprotocol waarmee verschillende elektronische
regeleenheden binnen een voertuig gegevens met elkaar kunnen uitwisselen
\autocite{CSSElectronicsCANBus2025}.

Een OBD-II-adapter vormt de verbinding tussen deze voertuiginterface en het
apparaat waarop de monitoringsoftware wordt uitgevoerd. Afhankelijk van het
type adapter kan de communicatie met de computer bijvoorbeeld via USB,
Bluetooth of Wi-Fi verlopen. Binnen deze bachelorproef wordt gebruikgemaakt
van een USB-adapter, zodat de OBD-II-communicatie rechtstreeks vanuit de
backend van de applicatie kan worden uitgevoerd.

\subsection{Parameter IDs}
Voor het opvragen van voertuiggegevens maakt OBD-II onder andere gebruik van
Parameter IDs, meestal afgekort als PIDs. Een PID identificeert een bepaalde
parameter die bij een elektronische regeleenheid kan worden opgevraagd.
Voorbeelden hiervan zijn het motortoerental, de voertuigsnelheid,
koelvloeistoftemperatuur en berekende motorbelasting
\autocite{CSSElectronicsOBD2025,mccord2011automotive}.

Bij het uitlezen van een parameter wordt een verzoek verstuurd waarin onder
andere wordt aangegeven welke PID moet worden opgevraagd. Het voertuig
beantwoordt dit verzoek met de beschikbare gegevens. De ontvangen bytes kunnen
vervolgens volgens de definitie van de betreffende PID worden omgezet naar een
bruikbare waarde \autocite{RFWirelessWorld2026}.

Niet alle PIDs zijn bij ieder voertuig beschikbaar. De ondersteuning is
afhankelijk van onder andere de fabrikant, het voertuigmodel, de motor en de
aanwezige elektronische regeleenheden. Hierdoor kan een applicatie niet zonder
meer aannemen dat iedere mogelijke OBD-II-parameter bij elk voertuig kan worden
uitgelezen \autocite{RIMPAS202055,Pmc2025}.

Voor een generieke monitoringapplicatie is het daarom belangrijk om eerst na te
gaan welke parameters daadwerkelijk door het aangesloten voertuig worden
ondersteund. In deze bachelorproef wordt dit onderzocht door meerdere
voertuigen uit te lezen en de beschikbare parameters met elkaar te vergelijken.
Op basis daarvan kan worden bepaald welke parameters geschikt zijn om verder in
de applicatie te gebruiken.

\subsection{Uitlezen van OBD-II-data}
OBD-II-communicatie verloopt in essentie volgens een request-responseprincipe.
De applicatie vraagt een bepaalde parameter op, waarna het voertuig een antwoord
met de overeenkomstige waarde terugstuurt. Wanneer meerdere parameters nodig
zijn, moeten bijgevolg meerdere gegevens bij het voertuig worden opgevraagd.

Het continu uitlezen van voertuigdata verschilt hierdoor van een systeem
waarbij sensoren automatisch op een vaste frequentie alle beschikbare waarden
naar de applicatie sturen. De snelheid waarmee nieuwe gegevens beschikbaar
komen, is onder andere afhankelijk van de gebruikte adapter, het
communicatieprotocol, het voertuig en het aantal parameters dat wordt
opgevraagd.

Onderzoek naar OBD-II-gebaseerde monitoringsystemen toont aan dat de interface
kan worden gebruikt voor het verzamelen en verwerken van voertuigdata tijdens
de werking van een voertuig \autocite{SageJournals2013,RIMPAS202055}. De
verkregen gegevens kunnen vervolgens door software worden verwerkt voor onder
andere monitoring, analyse en diagnose.

Deze eigenschappen zijn belangrijk binnen de context van deze bachelorproef.
De snelheid waarmee gegevens uit het voertuig kunnen worden verkregen, vormt
namelijk een eerste beperking voor de frequentie waarmee nieuwe informatie
uiteindelijk in het dashboard kan worden weergegeven.

\section{Realtime verwerking van voertuigdata}
Bij een realtime monitoringapplicatie worden veranderende gegevens regelmatig
verzameld en zo snel mogelijk aan de gebruiker weergegeven. In de context van
voertuigmonitoring betekent dit dat nieuwe OBD-II-waarden vanuit het voertuig
worden uitgelezen, door de backend worden verwerkt, naar de frontend worden
doorgestuurd en daar worden weergegeven.

De volledige datastroom kan vereenvoudigd worden voorgesteld als:

\begin{verbatim}
Voertuig -> OBD-II-adapter -> Backend -> Frontend -> Dashboard
\end{verbatim}

Iedere stap kan invloed hebben op hoe snel een nieuwe meetwaarde uiteindelijk
zichtbaar wordt. Niet alleen het uitlezen van de gegevens is daarom relevant,
maar ook de manier waarop de backend en frontend met elkaar communiceren.

Een belangrijke keuze hierbij is de updatefrequentie. Wanneer gegevens vaker
worden doorgestuurd, kan het dashboard veranderingen sneller weergeven. Een
hogere updatefrequentie betekent echter ook dat binnen dezelfde tijd meer
gegevens en berichten moeten worden verwerkt. Er bestaat daardoor een afweging
tussen de actualiteit van de weergegeven informatie en de hoeveelheid
communicatie die hiervoor nodig is.

Binnen deze bachelorproef worden REST API en WebSockets onderzocht als twee
mogelijke communicatiemethoden voor deze datastroom.

\section{REST API}
Representational State Transfer (REST) is een architecturale stijl voor
gedistribueerde systemen die werd beschreven door \textcite{Fielding2000}. 
REST wordt veel toegepast bij web-API's waarbij een client via HTTP communiceert met een server.

HTTP werkt volgens een request-responsemodel. Een client verstuurt een request
naar een server, waarna de server het verzoek verwerkt en een response
terugstuurt \autocite{RFC9110}. Bij een REST API worden resources doorgaans via
endpoints beschikbaar gesteld en kunnen HTTP-methoden zoals GET, POST, PUT en
DELETE worden gebruikt om met deze resources te werken.

Voor het opvragen van voertuigdata kan de frontend bijvoorbeeld periodiek een
GET-request naar een endpoint van de backend sturen. De backend leest of
verwerkt vervolgens de actuele voertuiggegevens en stuurt deze in de response
terug naar de frontend. De data kan daarbij bijvoorbeeld in JSON-formaat worden
doorgestuurd.

Wanneer REST wordt gebruikt voor gegevens die regelmatig veranderen, kan
polling worden toegepast. Hierbij verstuurt de client op een vast interval
telkens opnieuw een request naar de server. Bij een interval van één seconde
wordt bijvoorbeeld iedere seconde een nieuw request uitgevoerd.

Deze aanpak is relatief eenvoudig te implementeren en sluit goed aan bij het
gebruikelijke request-responsemodel van webapplicaties. Voor gegevens die zeer
frequent moeten worden bijgewerkt, betekent polling echter dat voortdurend
nieuwe HTTP-requests en responses nodig zijn. Naast de eigenlijke voertuigdata
bevat deze communicatie ook HTTP-informatie, waardoor extra netwerkverkeer
ontstaat.

Naarmate het pollinginterval kleiner wordt, neemt het aantal requests per
tijdseenheid toe. Hierdoor is het relevant om te onderzoeken welke invloed de
gekozen updatefrequentie heeft op onder andere de latency en het
netwerkverbruik van de toepassing.

\section{WebSockets}
WebSockets bieden een andere manier om communicatie tussen een client en server
te organiseren. Het WebSocket-protocol is gestandaardiseerd in RFC 6455 en
maakt bidirectionele communicatie mogelijk via een blijvende verbinding
\autocite{RFC6455}.

Bij het opzetten van een WebSocket-verbinding wordt eerst een handshake
uitgevoerd. Nadat de verbinding tot stand is gekomen, blijft deze actief en
kunnen zowel client als server berichten over dezelfde verbinding versturen.
Hierdoor hoeft niet voor iedere nieuwe dataset een afzonderlijke
HTTP-request-responsecyclus te worden uitgevoerd.

Voor een realtime monitoringapplicatie betekent dit dat de backend nieuwe
voertuigdata rechtstreeks via de bestaande WebSocket-verbinding naar de
frontend kan sturen. De frontend luistert naar binnenkomende berichten en kan
de ontvangen gegevens vervolgens onmiddellijk verwerken en weergeven.

De communicatie kan vereenvoudigd als volgt worden voorgesteld:

\begin{verbatim}
Frontend                       Backend
   |                              |
   | ----- WebSocket setup -----> |
   |                              |
   | <------ voertuigdata ------- |
   | <------ voertuigdata ------- |
   | <------ voertuigdata ------- |
   |             ...              |
\end{verbatim}

Het gebruik van een blijvende verbinding maakt WebSockets geschikt voor
toepassingen waarin regelmatig nieuwe gegevens tussen client en server moeten
worden uitgewisseld. Daartegenover staat dat de verbinding gedurende de
communicatie actief moet worden gehouden en dat de applicatie rekening moet
houden met het openen, sluiten en eventueel wegvallen van deze verbinding.

Binnen de ontwikkelde monitoringapplicatie is vooral relevant dat opeenvolgende
voertuiggegevens via dezelfde verbinding kunnen worden verzonden. Dit verschilt
van REST-polling, waarbij voor iedere update een nieuw HTTP-request wordt
uitgevoerd.

\section{REST API en WebSockets voor realtime monitoring}
REST API en WebSockets verschillen voornamelijk in de manier waarop de
communicatie tussen frontend en backend wordt georganiseerd. Bij REST-polling
neemt de frontend het initiatief om op vaste tijdstippen nieuwe gegevens op te
vragen. Bij WebSockets wordt een blijvende verbinding gebruikt waarover de
backend opeenvolgende gegevens naar de frontend kan sturen.

Beide methoden kunnen binnen een monitoringapplicatie worden gebruikt, maar hun
eigenschappen kunnen verschillen wanneer de updatefrequentie wordt verhoogd.
Bij REST neemt bij een kleiner pollinginterval het aantal HTTP-requests en
responses toe. Bij WebSockets neemt eveneens het aantal verzonden berichten
toe, maar deze berichten worden over een reeds bestaande verbinding
verstuurd.

Voor realtime voertuigmonitoring zijn voornamelijk latency en netwerkverbruik
relevant. Latency bepaalt hoeveel vertraging aanwezig is voordat nieuwe
informatie bij de frontend beschikbaar is. Een lage latency zorgt ervoor dat
het dashboard de actuele toestand van het voertuig zo snel mogelijk kan
weergeven.

Het gemiddelde alleen geeft echter niet altijd een volledig beeld van de
vertraging. Bij systemen waarin regelmatig gegevens worden verwerkt, kunnen
sporadische tragere resultaten belangrijk zijn. Percentielen kunnen worden
gebruikt om dergelijke uitschieters beter zichtbaar te maken
\autocite{Dean2013}. Binnen deze bachelorproef wordt daarom naast de gemiddelde
latency ook de P95-latency gebruikt.

Daarnaast is het netwerkverbruik relevant. Wanneer dezelfde voertuiggegevens
zeer frequent worden verstuurd, kan de bijkomende protocolinformatie een
groter aandeel van het totale netwerkverkeer vormen. Het verschil in
communicatiemodel tussen REST en WebSockets kan daardoor vooral bij kleinere
update-intervallen zichtbaar worden.

Op basis van de theoretische werking kan worden verwacht dat WebSockets bij
frequente updates minder protocoloverhead veroorzaken doordat niet voor iedere
dataset een afzonderlijk HTTP-request nodig is. Dit vormt echter een
theoretische verwachting en geen resultaat van dit onderzoek. Daarom worden
beide technieken later binnen dezelfde applicatie experimenteel met elkaar
vergeleken.

\section{Simulatie van voertuigdata}
Tijdens de ontwikkeling van een OBD-II-applicatie is het niet altijd praktisch
mogelijk om voortdurend een fysiek voertuig beschikbaar te hebben. Voor het
uitlezen van echte voertuigdata moeten een voertuig, OBD-II-adapter en computer
tegelijk beschikbaar zijn. Dit kan het ontwikkelen en testen van afzonderlijke
onderdelen van de applicatie bemoeilijken.

Een simulatieomgeving kan hiervoor een oplossing bieden. In plaats van
voortdurend gegevens uit een voertuig uit te lezen, worden voertuiggegevens
door software gegenereerd. De rest van de applicatie kan deze gegevens
vervolgens op een gelijkaardige manier verwerken als gegevens afkomstig van
een fysieke OBD-II-verbinding.

Een belangrijk voordeel hiervan is dat tests eenvoudiger kunnen worden
herhaald. Daarnaast kunnen bepaalde situaties gecontroleerd worden nagebootst
zonder dat daarvoor daadwerkelijk met een voertuig moet worden gereden. Voor
de ontwikkeling van de frontend kan hierdoor bijvoorbeeld worden gecontroleerd
hoe het dashboard reageert wanneer waarden veranderen.

Een simulator blijft echter een vereenvoudiging van de werkelijkheid.
Gesimuleerde gegevens bevatten niet noodzakelijk alle variaties, vertragingen
en onverwachte situaties die bij een echte OBD-II-verbinding kunnen optreden.
Daarom kan simulatie echte voertuigdata niet volledig vervangen.

Binnen deze bachelorproef wordt de simulator dan ook voornamelijk gebruikt als
hulpmiddel tijdens de ontwikkeling en functionele testing van de applicatie.
De uiteindelijke vergelijking tussen REST API en WebSockets wordt uitgevoerd
met echte voertuigdata, zodat ook de invloed van de werkelijke
OBD-II-databron aanwezig blijft in de testopstelling.

\section{Visualisatie van realtime voertuigdata}
Na het verzamelen en doorsturen van de voertuigdata moeten de gegevens op een
begrijpelijke manier aan de gebruiker worden gepresenteerd. Een dashboard kan
hiervoor actuele voertuigwaarden combineren met visuele elementen die een
snel overzicht van de toestand van het voertuig geven.

Bij realtime visualisatie worden de weergegeven waarden aangepast wanneer
nieuwe gegevens beschikbaar komen. Dit stelt andere eisen dan een statische
weergave, omdat de interface regelmatig opnieuw nieuwe informatie moet
verwerken.

Niet iedere parameter vereist daarbij dezelfde visualisatie. Een actuele
numerieke waarde kan rechtstreeks worden weergegeven, terwijl voor gegevens
waarbij veranderingen belangrijk zijn een grafische voorstelling nuttig kan
zijn. De keuze van de visualisatie hangt daarom af van het soort informatie dat
aan de gebruiker wordt getoond.

Binnen deze bachelorproef wordt een webgebaseerd dashboard ontwikkeld waarin
de uitgelezen OBD-II-gegevens realtime worden weergegeven. De frontend wordt
daarbij bijgewerkt met gegevens die via REST API of WebSockets vanuit de
backend worden ontvangen. De concrete implementatie van het dashboard wordt
later besproken in
Hoofdstuk~\ref{ch:ontwikkeling-van-een-realtime-monitoringdashboard}.

\section{Vergelijkende studie van bestaande onderzoeken}
In de literatuur zijn verschillende toepassingen terug te vinden waarin
OBD-II-data wordt gebruikt voor voertuigmonitoring, diagnose en analyse.
Hoewel deze onderzoeken raakvlakken hebben met deze bachelorproef, verschillen
ze in doelstelling, gebruikte technologieën en onderzoeksfocus.

Rimpas et al. onderzochten het gebruik van OBD-II-sensordata voor het monitoren
van de werking en het verbruik van voertuigen. Daarbij werden verschillende
OBD-II-parameters verzameld en geanalyseerd. Het onderzoek toont aan dat
OBD-II-data bruikbaar is voor het opvolgen van voertuiggedrag, maar richt zich
voornamelijk op de verzameling en analyse van de voertuiggegevens zelf
\autocite{RIMPAS202055}.

Yang et al. ontwikkelden een monitorings- en diagnosesysteem waarbij
voertuigdata via OBD-II en CAN werd verzameld en vervolgens op een
Android-smartphone werd weergegeven. Hiermee werd aangetoond dat voertuigdata
via OBD-II kan worden doorgestuurd naar een externe toepassing voor realtime
monitoring \autocite{SageJournals2013}.

Rybitskyi et al. onderzochten eveneens het gebruik van OBD-II binnen een
informatiesysteem voor voertuigdiagnose. Hierbij ligt de nadruk voornamelijk op
het verzamelen en gebruiken van diagnostische voertuiggegevens binnen een
informatiesysteem \autocite{Rybitskyi2023}.

Michailidis et al. geven in hun overzichtsstudie een bredere kijk op recente
toepassingen van OBD-II-data. Hieruit blijkt dat OBD-II-data onder andere wordt
gebruikt voor voertuigmonitoring, analyse van rijgedrag, brandstofoptimalisatie,
foutdetectie en machine-learningtoepassingen. De studie toont daarmee aan dat
OBD-II een bruikbare databron vormt voor verschillende datagedreven
voertuigtoepassingen \autocite{Pmc2025}.

Hoewel deze onderzoeken aantonen dat OBD-II-data kan worden verzameld,
verwerkt en gevisualiseerd, ligt hun onderzoeksfocus niet specifiek op de
vergelijking van verschillende communicatiemethoden tussen de backend en
frontend van een webgebaseerde realtime monitoringapplicatie. Vooral de invloed
van de communicatiemethode en de updatefrequentie op latency en netwerkverbruik
wordt binnen deze context minder rechtstreeks behandeld.

Tabel~\ref{tab:vergelijkende-studie} geeft een overzicht van de belangrijkste
verschillen tussen de besproken onderzoeken en deze bachelorproef.

\begin{table}[H]
    \centering
    \caption{Vergelijking van relevante onderzoeken met deze bachelorproef}
    \label{tab:vergelijkende-studie}
    \resizebox{\textwidth}{!}{
    \begin{tabular}{lccccc}
        \hline
        \textbf{Onderzoek} &
        \textbf{OBD-II} &
        \textbf{Realtime} &
        \textbf{Visualisatie} &
        \textbf{REST/WS vergelijking} &
        \textbf{Simulatie} \\
        \hline
        Rimpas et al. \autocite{RIMPAS202055}
            & Ja & Ja & Beperkt & Nee & Nee \\
        Yang et al. \autocite{SageJournals2013}
            & Ja & Ja & Ja & Nee & Nee \\
        Rybitskyi et al. \autocite{Rybitskyi2023}
            & Ja & Ja & Ja & Nee & Nee \\
        Michailidis et al. \autocite{Pmc2025}
            & Ja & Verschillend & Verschillend & Nee & Verschillend \\
        Deze bachelorproef
            & Ja & Ja & Ja & Ja & Ja \\
        \hline
    \end{tabular}
    }
\end{table}

De vergelijking toont aan dat de afzonderlijke onderdelen van deze
bachelorproef reeds in verschillende vormen in bestaande onderzoeken
voorkomen. Het uitlezen van OBD-II-data en het realtime weergeven van
voertuiggegevens zijn op zichzelf dus geen nieuwe concepten. Het verschil zit
voornamelijk in de combinatie van deze onderdelen en de specifieke vergelijking
van REST API en WebSockets binnen één OBD-II-monitoringapplicatie.

\section{Aanpak binnen deze bachelorproef}
Op basis van de bestaande literatuur wordt binnen deze bachelorproef gekozen
voor een praktische en experimentele aanpak. In plaats van uitsluitend een
monitoringapplicatie te ontwikkelen, worden verschillende onderdelen van de
volledige datastroom onderzocht.

Eerst wordt echte OBD-II-data verzameld bij meerdere voertuigen. Hierbij wordt
nagegaan welke parameters door de verschillende voertuigen worden ondersteund.
Dit is belangrijk omdat de beschikbare OBD-II-parameters afhankelijk kunnen
zijn van het voertuig. Op basis van deze vergelijking worden de gegevens
geselecteerd die verder binnen de applicatie worden gebruikt.

Vervolgens wordt een webgebaseerde monitoringapplicatie ontwikkeld. De
backend staat in voor de communicatie met de OBD-II-adapter en het verwerken
van de voertuigdata. De frontend ontvangt deze gegevens en geeft ze realtime
weer in een dashboard.

Daarnaast wordt een simulatieomgeving voorzien. Deze maakt het mogelijk om
tijdens de ontwikkeling voertuigdata te genereren wanneer geen fysiek voertuig
beschikbaar is. De simulator wordt voornamelijk gebruikt als hulpmiddel voor
ontwikkeling en functionele testing en vervangt de echte OBD-II-data niet bij
de uiteindelijke vergelijking van de communicatiemethoden.

Een belangrijk onderdeel van de aanpak is dat zowel REST API als WebSockets
binnen dezelfde applicatie worden geïmplementeerd. Hierdoor kunnen beide
communicatiemethoden worden getest terwijl de overige onderdelen van de
testopstelling zoveel mogelijk gelijk blijven.

De vergelijking wordt uitgevoerd met echte OBD-II-data. REST en WebSockets
worden getest met update-intervallen van 100, 250 en 1000 milliseconden. Voor
iedere configuratie worden latency en netwerkverbruik gemeten. Iedere test
wordt gedurende 60 seconden uitgevoerd en meerdere keren herhaald om de invloed
van toevallige schommelingen te beperken.

Door beide communicatiemethoden binnen dezelfde applicatie, met dezelfde
databron en onder vergelijkbare omstandigheden te testen, kan specifiek worden
onderzocht welke invloed de gekozen communicatiemethode en updatefrequentie
hebben op het doorsturen van realtime voertuigdata.

\section{Bijdrage van het onderzoek}
De bijdrage van deze bachelorproef ligt voornamelijk in het samenbrengen en
praktisch evalueren van verschillende aspecten van realtime
OBD-II-monitoring binnen één webapplicatie.

Ten eerste wordt onderzocht in welke mate de beschikbare OBD-II-parameters
verschillen tussen meerdere voertuigen. Hierdoor wordt in de praktijk zichtbaar
dat een applicatie voor voertuigmonitoring rekening moet houden met verschillen
in PID-ondersteuning.

Ten tweede wordt een volledige datastroom gerealiseerd waarbij voertuigdata
vanuit een fysieke OBD-II-verbinding via een backend naar een webgebaseerd
dashboard wordt doorgestuurd en realtime wordt weergegeven. Daarnaast wordt
een simulatieomgeving ontwikkeld die het mogelijk maakt om delen van de
applicatie zonder fysieke voertuigverbinding te ontwikkelen en te testen.

De belangrijkste experimentele bijdrage is de vergelijking van REST API en
WebSockets binnen dezelfde OBD-II-monitoringapplicatie. Beide technieken
worden onder vergelijkbare omstandigheden getest met verschillende
updatefrequenties. Daarbij worden latency en netwerkverbruik gemeten.

Het onderzoek heeft niet als doel om algemeen te bewijzen dat één van beide
communicatietechnieken in iedere realtime toepassing beter is. De resultaten
zijn specifiek verbonden aan de ontwikkelde applicatie en testopstelling. Ze
kunnen wel inzicht geven in de gevolgen van beide communicatieprincipes bij het
frequent doorsturen van OBD-II-data en ondersteunen zo de keuze van een
geschikte communicatiemethode voor gelijkaardige realtime
monitoringtoepassingen.